\section{Closing the force law}
\label{sec:closure}

The force law $\dot p=\kappa G^{\!\top}p$ carries its rate $\kappa$ as a constant put in by hand. A force law of this shape is familiar from
rigid bodies and ideal fluids, where the equation of motion is a geodesic equation on a Lie group and the rate is itself dynamical. This
section asks whether the interior law can be closed in that sense: whether there is a Lie--Poisson system whose bracket and inertia are built
from the same data that build $G$, and whose solutions contain the force law with $\kappa$ as a dynamical variable. A bracket or an inertia is
called \emph{derived from the Clifford structure} when it is equivariant under the algebra $\lG$ of \cref{thm:GLorentz} and built from
$\eta$, $G$ and the bivectors. The answers, in brief: no closure lives on the frame alone; the closure on one extra direction is unique; its
inertia is constrained by two requirements that exclude each other; and the dynamics it generates is a pendulum in the duality angle that
drifts the frame except on null data.

Throughout, $V=\R^{3,3}$, a covector $p$ is a column with $\langle p,X\rangle=p^{\!\top}X$, $\pi=\eta^{-1}p$, and $z(X)=x+iy$ for
$X=\sum_j(x_je_j+y_je_{j+3})$, so that $G$ acts as $i$. For a covector put
\begin{equation}
\label{eq:wdef}
C(p)=p^{\!\top}\eta^{-1}p,\qquad D(p)=p^{\!\top}G\eta^{-1}p,\qquad w(p):=q\bigl(z(\pi)\bigr)=C(p)-iD(p).
\end{equation}
In the field-strength reading, with $(\bB;\bE)=\pi$, $w=|\bB|^2-|\bE|^2+2i\,\bB\cdot\bE$ is the complex invariant of the field carried by the
momentum.

\subsection{The force law as a dual transport}

Euler--Arnold equations are geodesic equations on a Lie group with a one-sided invariant metric \cite{Arnold66,ArnoldKhesin98}; for a
left-invariant metric they read $\dot\mu=\ad_w^{\!\top}\mu$ with $\mu=Aw$ for an inertia $A$. The rigid body is the case of $SO(3)$, where this is
$\dot{\mathbf M}=\mathbf M\times\boldsymbol\Omega$, and incompressible Euler flow the case of the volume-preserving diffeomorphisms.

\begin{proposition}[The dictionary]
\label{prop:EAdict}
\begin{enumerate}[label=(\roman*),leftmargin=2.2em]
\item Since $\langle G^{\!\top}p,v\rangle=\langle p,Gv\rangle$, the law $\dot p=\kappa G^{\!\top}p$ transports the covector $p$ by the transpose of $\kappa G$.
It is the Lie--Poisson equation $\dot p=\ad_w^{\!\top}p$ of a Lie bracket on the frame exactly when $\kappa G=\ad_w$ is an inner derivation.
\item For a bracket with trivial centre on the $6n$-dimensional frame the inner derivations fill a $6n$-dimensional subspace of $\so{3n,3n}$, of
dimension $3n(6n-1)$, when $\eta$ is invariant: a fraction $2/(6n-1)$.
\item When $\kappa G=\ad_u$ and $p=Au$, the shift operators $\alpha_\pm=\partial_\tau\pm\kappa G$ are the constrained variations
$\delta u=\dot\xi\pm[u,\xi]$ of Euler--Poincar\'e reduction \cite{MarsdenRatiu99,HolmMarsdenRatiu98}, in the body and in the spatial frame.
\end{enumerate}
\end{proposition}

\begin{proposition}[The binding root is inner for no bracket]
\label{prop:EAnotinner}
For every Lie bracket and every $w$, $\ad_w(w)=0$, so every inner derivation is singular. The binding root satisfies $G^2=-\id$ and is
invertible, so $\kappa G$ is inner for no Lie bracket on $\R^{3n,3n}$, invariant or not, and the force law with that $G$ is never a
Lie--Poisson equation on the frame itself. It is one on the extension $\R\hat g\ltimes\R^{3n,3n}$, the frame taken abelian and $\hat g$ acting on
it by $G$, where it is the Lie--Poisson equation of the fixed element $\kappa\hat g$ restricted to the ideal.
\end{proposition}

\begin{theorem}[Transport and factorisation are complementary]
\label{thm:EAdichotomy}
$\End(\R^{3n,3n})=\so{3n,3n}\oplus\{G:G=G^\dagger\}$, of dimensions $3n(6n-1)+3n(6n+1)=36n^2$. Along $\dot p=\kappa G^{\!\top}p$ the Casimir
$C=\eta^{-1}(p,p)$ changes at the rate $dC/d\tau=2\kappa\,p^{\!\top}G\eta^{-1}p$, which vanishes for every $p$ exactly when $G$ lies in the first
summand, while factorisation in the split form requires the second. The binding root lies in the second summand and moves the Casimir at the rate
$dC/d\tau=-4\kappa\,\bB\cdot\bE$ in the variables $(\bB;\bE)=\eta p$.
\end{theorem}

At the first rung the frame itself is a Lie algebra. With $L_i=e_i$ and $K_i=e_{i+3}$ obeying the brackets of \cref{thm:emfield}, $\R^{3,3}$
becomes $\so{1,3}$ with Killing form $-4\eta$; in the variable $z=\bB+i\bE$ the bracket is the complex cross product, $G$ is multiplication by $i$
and lies in the centroid, $[Gx,y]=G[x,y]$, and $z\cdot z$ is the form $q$. The inner derivations fill six of the fifteen dimensions of
$\so{3,3}$. Closing the force law into an Euler--Arnold equation on the frame would need a bracket for which the shift is inner, and by
\cref{prop:EAnotinner} there is none. A closure must leave the frame.

\subsection{No closure on the frame; the rotor is unique}

\begin{proposition}[Equivariant maps and forms]
\label{prop:lGhoms}
For $\lG=\operatorname{span}\{L_i,K_i\}$: $\Hom_{\lG}(V,V)=\operatorname{span}\{\id,G\}$ and $\Hom_{\lG}(V,\R)=0$;
$\Hom_{\lG}(\Lambda^2V,V)=\operatorname{span}\{\beta_0,G\beta_0\}$, where $\beta_0(X,Y)$ is the complex cross product $z(X)\times z(Y)$ returned to
$V$, and $\Hom_{\lG}(\Lambda^2V,\R)=0$; and the $\lG$-invariant symmetric forms on $V^*$ are $\operatorname{span}\{\eta^{-1},G\eta^{-1}\}$.
\end{proposition}

\begin{theorem}[The brackets on the frame]
\label{thm:brackets}
The $\lG$-equivariant Lie brackets on $V$ are $\beta_L=L\beta_0$ with $L=a+bG$. The bracket $\beta_0$ makes $(V,\beta_0)\cong\lG$, its derivations
are the six inner ones, its centroid is $\operatorname{span}\{\id,G\}\cong\C$, so $G$ is not a derivation, and $\beta_L$ has Killing form
$-4\bigl[(a^2-b^2)\eta+2ab\,\eta G\bigr]$.
\end{theorem}

\begin{proof}
By \cref{prop:lGhoms} every equivariant antisymmetric bilinear map is $(c_1+c_2G)\beta_0$. Since $L$ lies in the centroid,
$\beta_L(\beta_L(X,Y),Z)=L^2\beta_0(\beta_0(X,Y),Z)$, and $\beta_L$ inherits the Jacobi identity.
\end{proof}

\begin{theorem}[No closure on the frame alone]
\label{thm:noframe}
Let $L=a+bG\neq0$. For every symmetric $N$, with $\mathfrak g^*$ identified with $\mathfrak g$ through $\eta$, the Euler--Arnold field
$\beta_0(M,NM)$ conserves $\eta(M,M)$ and $\eta(M,GM)$, whereas along the force law, in the form $\dot M=\kappa GM$,
$\tfrac{d}{d\tau}\eta(M,M)=2\kappa\,\eta(M,GM)$ and $\tfrac{d}{d\tau}\eta(M,GM)=-2\kappa\,\eta(M,M)$. No Hamiltonian flow on $(V,\beta_L)^*$ coincides
with the force law on an open set, and this remains true with a state-dependent rate.
\end{theorem}

\begin{proof}
The two quadratic forms are the Casimirs of $(V,\beta_L)^*$ and are conserved by every Hamiltonian flow, while the force law rotates them into
each other. Equality on an open set would make both vanish there, that is $z\cdot z=0$ on an open set, which no nonzero polynomial does. For a
state-dependent rate, the equations $\beta_0(M_s,NM_s)=f_s\,GM_s$ at twelve rational null vectors already force every $f_s=0$ and
$N\in\operatorname{span}\{\id,G\}$, for which $\beta_0(M,NM)\equiv0$.
\end{proof}

A closure therefore needs a direction beyond the frame. Adjoin one generator $\hat g$, the \emph{rotor}, on which $\lG$ acts trivially; the rate
$\kappa$ is a Lorentz scalar, since $G$ commutes with $\lG$.

\begin{theorem}[Uniqueness of the closure]
\label{thm:closureunique}
Let $\mathfrak s=\R\hat g\oplus V$ carry an $\lG$-equivariant Lie bracket for which $\ad_{\hat g}$ does not vanish on $V$. Then
$[\hat g,X]=(a+bG)X$ and $[X,Y]=0$ for $X,Y\in V$, with $(a,b)\neq(0,0)$. If moreover the rotor rotates the frame covector without dilating it,
as the force law does, then $a=0$ and, after rescaling $\hat g$, $\mathfrak s\cong\R\ltimes_GV$: $[\hat g,X]=GX$, $[X,Y]=0$.
\end{theorem}

\begin{proof}
Equivariance gives $[\hat g,X]=AX+\alpha(X)\hat g$ with $A\in\operatorname{span}\{\id,G\}$ and $\alpha=0$, and $[X,Y]=(c_1+c_2G)\beta_0(X,Y)$, the
$\hat g$-component vanishing by \cref{prop:lGhoms}. The Jacobi identity for $(\hat g,X,Y)$ reads $(a+bG)(c_1+c_2G)\beta_0(X,Y)=0$, a linear system in
$(c_1,c_2)$ with determinant $a^2+b^2\neq0$, so $c_1=c_2=0$. On $\mathfrak s^*$ the rotor acts on $p$ through $(a+bG)^{\!\top}$, and
$\tfrac{d}{d\tau}|p|^2=2a\,(h/I_0)|p|^2$ because $G$ is antisymmetric; a pure rotation needs $a=0$.
\end{proof}

The extension of \cref{prop:EAnotinner} is therefore not one choice among several: it is the only equivariant closure. The frame is abelian,
all the dynamics sits in the rotor, and the rotor acts by the binding root itself. With several constituents the argument runs factor by
factor, since equivariant maps between distinct constituents vanish, and a common rotation rate gives $\mathfrak s_n=\R\ltimes_{G_n}V_n$ with
$G_n=\bigoplus_kG$.

\subsection{The Lie--Poisson dynamics, and the two interiors as clamped reductions}

\begin{proposition}[Lie--Poisson equations]
\label{prop:LP}
On $\mathfrak s^*$ with coordinates $(h,p)$ let $H=h^2/(2I_0)+\tfrac12p^{\!\top}Np$, with $I_0>0$ the rotor inertia and $N$ a symmetric form on
$V^*$, the inverse frame inertia. The Lie--Poisson equations are
\begin{equation}
\label{eq:LProtor}
\dot p=\frac h{I_0}\,G^{\!\top}p,\qquad \dot h=-\,p^{\!\top}GNp ,
\end{equation}
so the force law holds with the dynamical rate $\kappa=h/I_0$, and $|p|^2$ is a Casimir. For $p\neq0$ the coadjoint orbit through $(h,p_0)$ is the
cylinder $\R\times\{e^{\theta G}p_0\}\cong T^*S^1$, on which $(\theta,h)$ are canonical for $h^2/(2I_0)+U(\theta)$ with
$U(\theta)=\tfrac12p(\theta)^{\!\top}Np(\theta)$ and $p(\theta)=e^{-\theta G}p_0$: every reduced motion is a motion of the rotor angle alone.
\end{proposition}

\begin{proposition}[The two interiors are clamped reductions]
\label{prop:groupred}
On the group $S$ with product $(\theta,x)(\theta',x')=(\theta+\theta',x+e^{\theta G}x')$ let $L=\tfrac12I_0\dot\theta^2+\tfrac m2\eta(u,u)$ with body velocity
$u=e^{-\theta G}\dot x$, and put $X=e^{-\theta G}x$. Then $u=\dot X+\dot\theta\,GX$, the Euler--Lagrange equations are \eqref{eq:LProtor} with $N=\eta^{-1}/m$,
and clamping $\theta=\kappa\tau$ gives
\[
L=\tfrac12I_0\kappa^2+\tfrac m2\,\eta\bigl(\dot X+\kappa GX,\ \dot X+\kappa GX\bigr),
\]
the split-form interior \eqref{eq:splitL}. With $\tfrac m2|u|^2$ in place of $\tfrac m2\eta(u,u)$ the clamped Lagrangian is the positive-form interior
\eqref{eq:alpha} with $g=G$.
\end{proposition}

So the two interiors of \cref{sec:shifted,sec:neutral} are the clamped reductions of one closure, with inverse inertias $N=\eta^{-1}/m$ and
$N=\id/m$. Which inertias the Clifford data permit is the next question.

\subsection{Two requirements on the inertia}

\begin{definition}[The two requirements]
\label{def:inertiareq}
A symmetric $N$ on $V^*$ satisfies the \emph{Lorentz requirement} if it is $\lG$-invariant, $rN+Nr^{\!\top}=0$ for all $r\in\lG$, and the
\emph{duality requirement} if the rate $h/I_0$ is conserved for all data. Write $N=N_c+N_a$ with $N_c=\tfrac12(N-GNG)$ commuting and
$N_a=\tfrac12(N+GNG)$ anticommuting with $G$.
\end{definition}

\begin{theorem}[The two requirements exclude each other]
\label{thm:inertiareq}
\begin{enumerate}[label=(\roman*),leftmargin=2.2em]
\item $N_c$ and $N_a$ are symmetric; the commuting forms have dimension $9$ and the anticommuting ones $12$.
\item The Lorentz requirement gives $N=(a\eta^{-1}+bG\eta^{-1})/m$, a two-dimensional family inside the anticommuting forms, with frame energy
$\tfrac12p^{\!\top}Np=\Re\bigl((a+ib)w(p)\bigr)/(2m)$.
\item The duality requirement holds exactly when $[N,G]=0$. Among the forms invariant under the compact part of $\lG$ the only ones commuting with
$G$ are the multiples of $\id$: the positive form.
\item The two requirements meet in $0$, and they do not exhaust the symmetric forms: $N=(2\,\id+\eta^{-1})/m$ is positive definite and satisfies
neither.
\end{enumerate}
\end{theorem}

In the field-strength reading the dichotomy is the classical one. The Lorentz class consists of the two invariants $|\bB|^2-|\bE|^2$ and
$\bB\cdot\bE$; the duality class contains the energy density $|\bB|^2+|\bE|^2$; and a duality rotation preserves the energy density while it
rotates the Lagrangian density \cite{Calkin65,DeserTeitelboim76,GaillardZumino81}. The split form lies in the Lorentz class and the positive form
in the duality class.

\begin{theorem}[The frame under the closure]
\label{thm:drift}
Reconstruct the group curve from the body velocity,
\[
(\theta,x)^{-1}\tfrac{d}{d\tau}(\theta,x)=(h/I_0,\,Np),\qquad \theta(0)=0 .
\]
Then $p=e^{-\theta G}p_0$ and
\[
\dot x=\bigl(N_c+e^{2\theta G}N_a\bigr)p_0,\qquad \dot h=-\,p_0^{\!\top}Ge^{2\theta G}N_ap_0 .
\]
With the rotor clamped at $\theta=\kappa\tau$, $\kappa\neq0$,
\[
x(\tau)=x_0+\tau\,N_cp_0+\bigl(e^{2\kappa\tau G}-\id\bigr)(2\kappa G)^{-1}N_ap_0 ,
\]
so the frame stays bounded for all data exactly when $N_c=0$ and drifts at the velocity $N_cp_0$ otherwise, while the rate is constant for all data
exactly when $N_a=0$.
\end{theorem}

A clamped rate that is also the dynamical one, and a bounded clamped frame, are therefore exclusive requirements: each is met by one summand of
$\mathrm{Sym}(V^*)=\mathrm{Sym}_c\oplus\mathrm{Sym}_a$. For $N=\id/m$ the rate is conserved and the frame drifts at $p_0/m$, which is the
non-recurrence of \cref{thm:virial}, $X=e^{-\kappa\tau G}(X_0+\tau p_0/m)$; for $N=\eta^{-1}/m$ the clamped frame is bounded, as in \cref{thm:neutral},
but the rate is conserved only on null data.

\subsection{The free rotor: a pendulum in the duality angle}

\begin{theorem}[The duality pendulum]
\label{thm:pendulum}
Let $N=(a\eta^{-1}+bG\eta^{-1})/m$ and write $(a+ib)w(p_0)=R_0e^{i\varphi_0}$, $R_0\ge0$. Along the flow $z(\pi)$ turns as $e^{i\theta}$ and
$w(p)=e^{2i\theta}w(p_0)$, the frame energy is $U(\theta)=\tfrac{R_0}{2m}\cos(2\theta+\varphi_0)$, and
\[
I_0\ddot\theta=\frac{R_0}m\sin(2\theta+\varphi_0),\qquad E_\theta=\tfrac12I_0\dot\theta^2+\frac{R_0}{2m}\cos(2\theta+\varphi_0)=H .
\]
With $\psi:=2\theta+\varphi_0$ this is the pendulum $\ddot\psi=(2R_0/mI_0)\sin\psi$, stable at $\psi=\pi$. The rotor angle is the duality angle of
the momentum field, and the potential is its Lorentz-invariant frame energy.
\end{theorem}

\begin{figure}[t]
\centering
\begin{tikzpicture}
\begin{axis}[width=12cm,height=6.4cm,xmin=-0.3,xmax=6.6,ymin=-3.0,ymax=3.0,
  axis lines=middle,xlabel={$\psi=2\theta+\varphi_0$},ylabel={$\dot\psi$},
  xtick={0,3.14159,6.28318},xticklabels={$0$,$\pi$,$2\pi$},ytick=\empty,
  xlabel style={anchor=north west},ylabel style={anchor=south west}]
\foreach \E/\a in {-0.5/2.0944,0/1.5708,0.5/1.0472}{
  \addplot[tcol,thick,domain=\a:{2*pi-\a},samples=160] {sqrt(max(0,2*(\E-cos(deg(x)))))};
  \addplot[tcol,thick,domain=\a:{2*pi-\a},samples=160] {-sqrt(max(0,2*(\E-cos(deg(x)))))};
}
\addplot[scol,very thick,domain=0:2*pi,samples=200] {2*abs(sin(deg(x/2)))};
\addplot[scol,very thick,domain=0:2*pi,samples=200] {-2*abs(sin(deg(x/2)))};
\addplot[gcol,thick,domain=0:2*pi,samples=160] {sqrt(2*(1.9-cos(deg(x))))};
\addplot[gcol,thick,domain=0:2*pi,samples=160] {-sqrt(2*(1.9-cos(deg(x))))};
\addplot[only marks,mark=*,mark size=2.2pt,black] coordinates {(3.14159,0)};
\addplot[only marks,mark=o,mark size=2.6pt,black] coordinates {(0,0) (6.28318,0)};
\end{axis}
\end{tikzpicture}
\caption{Phase portrait of the duality pendulum (\cref{thm:pendulum}), in units with $2R_0/(mI_0)=1$; the filled dot is the stable point
$\psi=\pi$ and the open dots the unstable one. Blue: librations, in which the rate $\kappa=\dot\theta$ changes sign periodically and the duality
rotation reverses. Red: the separatrix. Green: rotations, in which the rotor turns for all time with a modulated rate. The constant-rate law of \cref{sec:neutral} is the limit of large rotor energy, far above the
separatrix (\cref{cor:multiscale}).}
\label{fig:pendulum}
\end{figure}

\begin{corollary}[When the rate is constant]
\label{cor:constrate}
With Lorentz-class inertia the rate is constant and nonzero along the flow exactly when $R_0=0$ and $h(0)\neq0$; $R_0=0$ means $(a,b)=(0,0)$ or
$w(p_0)=0$, that is $|\bB|=|\bE|$ and $\bB\cdot\bE=0$ for the field of $\pi_0$: a null field, the class that contains plane waves
\cite{Synge56,Robinson61}.
\end{corollary}

\begin{corollary}[The multiscale criterion]
\label{cor:multiscale}
Let $R_0>0$, $\kappa_0:=\dot\theta(0)\neq0$, $\psi_0=2\theta(0)+\varphi_0$ and $\varepsilon:=R_0/(mI_0\kappa_0^2)$. Then
$\kappa^2=\kappa_0^2\bigl[1+\varepsilon(\cos\psi_0-\cos\psi)\bigr]$. The rotor turns for all time exactly when $\varepsilon(1-\cos\psi_0)<1$, and $\varepsilon<\tfrac12$
suffices for every phase, in which case $|\kappa^2/\kappa_0^2-1|\le2\varepsilon$. At $\varepsilon(1-\cos\psi_0)=1$ the motion is the separatrix, and beyond it
the rotor librates and the duality rotation reverses. The constant-rate law is the limit $\varepsilon\to0$, and the multiscale requirement is
\begin{equation}
\label{eq:multiscale}
I_0\,\kappa_0^2\ \gg\ R_0/m .
\end{equation}
\end{corollary}

A clamped rotor is therefore an idealisation of a heavy one. It is not a harmless idealisation, because the frame feels the modulation.

\begin{theorem}[The free rotor drifts the frame]
\label{thm:freedrift}
Let $N$ be of Lorentz class with $R_0>0$ and $\kappa_0\neq0$, in the rotating regime. During one period of the rate, while $\theta$ changes by $\pi$,
the frame is displaced by
\[
\Delta x=\lambda\,e^{-\varphi_0G}Np_0,\qquad \lambda=\tfrac12\int_0^{2\pi}\frac{\cos\psi\,d\psi}{|\kappa_0|\sqrt{1+\varepsilon(\cos\psi_0-\cos\psi)}}>0,
\]
and its mean velocity is $\tfrac\varepsilon4e^{-\varphi_0G}Np_0+O(\varepsilon^2)$. Every rotating solution with $R_0\neq0$ therefore has an unbounded frame.
\end{theorem}

\begin{proof}
For Lorentz-class $N$, $N=N_a$ and $\dot x=e^{2\theta G}Np_0$. In the rotating regime $\dot\theta$ never vanishes, so $d\tau=|d\theta|/|\dot\theta|$, and
with $e^{2\theta G}=e^{-\varphi_0G}e^{\psi G}$ and $e^{\psi G}=\cos\psi+G\sin\psi$ the $\sin\psi$ part integrates to zero because $1/|\dot\theta|$ depends on
$\cos\psi$ only. Since $1/|\dot\theta|$ is a strictly increasing function of $\cos\psi$ when $R_0>0$, the integrand of $\lambda$ exceeds its value at
$\cos\psi=0$ wherever $\cos\psi>0$ and falls short of it where $\cos\psi<0$, so $\lambda>0$. Expanding to first order gives
$\lambda=\pi\varepsilon/(4|\kappa_0|)+O(\varepsilon^2)$ over a period $\pi/|\kappa_0|+O(\varepsilon)$.
\end{proof}

\begin{example}[A drift in closed form]
\label{ex:drift}
Take $p_0=e^1$ and $m=I_0=\kappa_0=1$, so that $w=1$, $\varphi_0=0$, $\varepsilon=1$ and $\dot\theta^2=1+2\sin^2\theta$. Writing
$(1-2s)/\sqrt{1+2s}=2/\sqrt{1+2s}-\sqrt{1+2s}$ with $s=\sin^2\theta$, the displacement per turn of the rotor is
\[
\Delta x=\bigl(8K(-2)-4E(-2)\bigr)e_1\approx0.6336081022\,e_1 ,
\]
with $K(m)=\int_0^{\pi/2}(1-m\sin^2t)^{-1/2}dt$ and $E(m)=\int_0^{\pi/2}(1-m\sin^2t)^{1/2}dt$ the complete elliptic integrals in the parameter $m$
\cite{DLMF}. A numerical integration of the full system, frame included, reproduces this value to ten digits and returns $h$ and $p$ to their
initial values after each turn.
\end{example}

The clamped frame of \cref{thm:drift} moves on a circle of diameter $|Np_0|/|\kappa_0|$, while the free rotor adds $\tfrac{\pi\varepsilon}{2|\kappa_0|}|Np_0|$
per turn to first order, so the constant-rate picture tracks the bounded orbit for about $2/(\pi\varepsilon)$ turns. With Lorentz-class inertia and a
free rotor, a constant rate and a bounded frame coexist exactly on null data.

\begin{proposition}[Several constituents and one rotor]
\label{prop:several}
Let one rotor act on $n$ constituents with $N_k=(a_k\eta^{-1}+b_kG\eta^{-1})/m_k$ and $w_k=w(p^{(k)})$. Then $\dot w_k=2i(h/I_0)w_k$ for every $k$, all
forces carry the common phase $e^{2i\theta}$, and with $R_{\rm eff}e^{i\varphi_{\rm eff}}:=\sum_k(a_k+ib_k)w_k(0)/m_k$ the rotor obeys
$I_0\ddot\theta=R_{\rm eff}\sin(2\theta+\varphi_{\rm eff})$. The \emph{balance surface} $R_{\rm eff}=0$ is preserved by the flow; on it the rate is exactly
constant although no $w_k$ need vanish, and every frame moves on a bounded orbit. With equal inertias, $p^{(2)}=-Gp^{(1)}$ gives $w_2=-w_1$.
\end{proposition}

The Clifford data $\eta$, $G$, $G\eta^{-1}$ and $\beta_0$ are integer matrices. The rescaling $(I_0,m)\mapsto(sI_0,tm)$ preserves admissibility and
multiplies $\varepsilon$ by $1/(st)$, so no value of $\varepsilon$, and no ratio among $I_0$, $m$ and $|a+ib|$, is determined by the Clifford structure:
the form of the inertia is derived, and \eqref{eq:multiscale} is an input.

\subsection{Three constituents: the closure symmetry and the two projectors}

At $n=3$ write $e_A=e^{(k)}_i$ with $A=6(k-1)+i$; the temporal indices are $\{1,2,3,7,8,9,13,14,15\}$, each with its partner $i'=i+3$, and for
temporal $i<j$ put
\[
\mathrm{rot}(i,j):=e_{i'}e_{j'}-e_ie_j,\qquad \mathrm{bst}(i,j):=e_je_{i'}-e_ie_{j'} .
\]
For one constituent these span $\lG$, with $\mathrm{rot}(2,3)=2L_1$ and $\mathrm{bst}(2,3)=2K_1$ and cyclically up to sign. In this numbering the
weak directions of $P_*$ are $W=\{7,13,14,15\}$, the colour directions $\{10,11,12,16,17,18\}$, and $E$ denotes the span of the temporal directions
$\{7,8,9,13,14,15\}$ of constituents two and three; $G_3=G\oplus G\oplus G$.

\begin{theorem}[The closure symmetry is $\so{9,\C}$]
\label{thm:so9C}
$G_3$ is an anti-isometry of $\eta_3$ with $G_3^2=-\id$, and its centraliser in $\so{9,9}$ is $72$-dimensional, spanned by the $36$ paired rotations
and the $36$ paired boosts. In the basis of the nine temporal directions followed by their partners every element is
$\bigl(\begin{smallmatrix}R&S\\-S&R\end{smallmatrix}\bigr)$ with $R,S$ antisymmetric, and $R-iS$ identifies the centraliser with $\so{9,\C}$, whose
compact part consists of the paired rotations.
\end{theorem}

In the field-strength reading $V_3\cong\C^9$ carries three field strengths, and $\so{9,\C}$ consists of the complex orthogonal transformations of
$\sum_kz^{(k)}\cdot z^{(k)}$.

\begin{theorem}[No common Lorentz algebra]
\label{thm:nocommon}
The centraliser of $G_3$ meets the interior $\so{1,3}$ of the first constituent in $0$; the elements commuting with both $G_3$ and that algebra form
the $30$-dimensional $\so{6,\C}$ of paired elements of constituents two and three. Under $\lG^{(1)}$ the complexified spinor module is
$128\cdot(\tfrac12,\tfrac12)$, with no Weyl spinor.
\end{theorem}

No nonzero Lorentz generator leaves both the closure and the spinor structure of matter invariant: the two Lorentz algebras are different
subalgebras of $\so{9,9}$ with zero intersection, which is \cref{prop:notconj} lifted to three constituents. A Pati--Salam generator acts on the
closure variables only through its grade-two part, since $\ad_Z$ maps vectors to vectors exactly when $Z$ has grades zero and two. The intersection
with the Pati--Salam algebra can then be computed in general.

\begin{theorem}[The Pati--Salam algebra and the binding root]
\label{thm:PSbinding}
Let $U$ be any four-dimensional subspace of the span of the eight internal temporal directions, and $\mathfrak{ps}_U=\so U\oplus\so6_C$. Then neither
summand meets the centraliser of $G_3$, and
\[
\mathfrak{ps}_U\cap\Cent(G_3)=\bigl\{A+G_3AG_3^{-1}:\ A\in\so{U\cap E}\bigr\}\cong\so{U\cap E},
\]
of dimension $\binom d2\in\{1,3,6\}$ with $d=\dim(U\cap E)$. For the Pati--Salam algebra of $P_*$, $U=\operatorname{span}W\subset E$, and the
binding-compatible part is the diagonal $\so4_\Delta=\operatorname{span}\{\mathrm{rot}(i,j):i<j\in W\}\cong\su2\oplus\su2$, compact and non-abelian.
\end{theorem}

\begin{proof}
Let $A\in\so U$ and $B\in\so6_C$ with $[A+B,G_3]=0$, and split the internal temporal span as $E\oplus F$ with $F=\operatorname{span}\{e_2,e_3\}$. $G_3$ maps
$F$ into the interior plane $\operatorname{span}\{e_5,e_6\}$, where $A$ and $B$ vanish, so $Av=0$ for $v\in F$. For $v\in E$, $G_3v$ is a colour direction
and $Bv=0$, so $G_3Av=BG_3v$ lies in the colour span and $Av\in E$. Hence $A\in\so{U\cap E}$ and $B=G_3AG_3^{-1}$ on the colour span, and conversely.
Finally $d\ge4+6-8=2$.
\end{proof}

Neither Pati--Salam factor contains an element commuting with the binding root; the compatible part pairs every rotated temporal direction with
its spatial partner, which the colour $\su4$, acting on spatial directions alone, never does.

\begin{theorem}[What survives each projector]
\label{thm:survive}
\begin{enumerate}[label=(\roman*),leftmargin=2.2em]
\item The part of $\Cent(G_3)$ commuting with $\Pfam$ is $\operatorname{span}\{\mathrm{rot}(2,3),\mathrm{rot}(8,9),\mathrm{rot}(14,15),\mathrm{bst}(14,15)\}$,
abelian of dimension four. The kinetic condition holds for the two internal rotations and fails for $\mathrm{bst}(14,15)$, so the gauge-admissible
part is $\uu1^2=\operatorname{span}\{\mathrm{rot}(8,9),\mathrm{rot}(14,15)\}$.
\item The part commuting with $P_*$ is $\operatorname{span}\{\mathrm{rot}(2,3),\mathrm{rot}(8,9)\}\oplus\so4_\Delta$, compact of dimension eight, and its
internal part $\so4_\Delta\oplus\R\,\mathrm{rot}(8,9)$ acts faithfully and satisfies the kinetic condition.
\item Under either projector the frame Lorentz algebra $\lG^{(1)}$ keeps only $\mathrm{rot}(2,3)=e_5e_6-e_2e_3$, a rotation of the interior plane
$(e_5,e_6)$ locked to the rotation of the surplus pair.
\end{enumerate}
\end{theorem}

The rotor acts through $G_3$ and commutes with all of $\Cent(G_3)$, so a one-rotor closure is compatible with either projector. What differs is the
symmetry each leaves: an abelian one for $\Pfam$, and for $P_*$ the non-abelian diagonal of weak isospin and a colour $\so4$.

\begin{table}[t]
\centering
\small
\begin{tabular}{@{}lccccccccc@{}}
\toprule
$(q_1,q_2)$ & $(-1,-1)$ & $(-1,0)$ & $(-1,1)$ & $(0,-1)$ & $(0,0)$ & $(0,1)$ & $(1,-1)$ & $(1,0)$ & $(1,1)$\\
\midrule
left, $\Gamma=-1$ & 1 & 0 & 1 & 1 & 2 & 1 & 0 & 2 & 0\\
right, $\Gamma=+1$ & 0 & 2 & 0 & 1 & 2 & 1 & 1 & 0 & 1\\
\bottomrule
\end{tabular}
\caption{Weyl content of the block of $\Pfam$ under the two surviving charges $Q_1=-\tfrac12J\,\mathrm{rot}(8,9)$ and $Q_2=-\tfrac12J\,\mathrm{rot}(14,15)$:
the number of Weyl components with charges $(q_1,q_2)$ on each chirality half.}
\label{tab:closurecontent}
\end{table}

\begin{theorem}[Chirality and anomalies of what survives]
\label{thm:closurechiral}
With anomaly coefficients $\mathcal A(f)=\tfrac14\bigl[\Tr(fP_L)-\Tr(fP_R)\bigr]$ and charges $-JB$:
\begin{enumerate}[label=(\roman*),leftmargin=2.2em]
\item On $\Pfam$ the content is \cref{tab:closurecontent}; $\mathcal A\bigl((\alpha Q_1+\beta Q_2)^3\bigr)=-12\alpha\beta^2$ and $\mathcal A(\alpha Q_1+\beta Q_2)=0$, and
the left and right charge multisets coincide exactly when $\alpha\beta=0$. The anomaly-free subalgebras of $\uu1^2$ are $0$, $\R Q_1$ and $\R Q_2$,
each vector-like.
\item On $P_*$, with $\so4_\Delta=\su2_+\oplus\su2_-$ and $Q=-\tfrac12J\,\mathrm{rot}(8,9)$, the halves are
$L=(0,0)_0\oplus(\tfrac12,\tfrac12)_{+1}\oplus(1,0)_0$ and $R=(0,0)_0\oplus(\tfrac12,\tfrac12)_{-1}\oplus(0,1)_0$. For $H=\alpha H^+_3+\beta H^-_3+\gamma Q$,
$\mathcal A(H^3)=2\gamma(3\alpha^2+3\beta^2+4\gamma^2)$ and $\mathcal A(H)=8\gamma$, so the anomaly-free subalgebras are exactly those of $\so4_\Delta$, which
has neither a perturbative nor a global anomaly.
\item The two halves of $P_*$ carry inequivalent representations of $\so4_\Delta$, so it acts chirally; but every summand is a real representation
of $\su2\oplus\su2$, so the left-handed content is real and gauge-invariant Majorana masses exist \cite{Suzuki2000}. The chirality of $\so4_\Delta$
on $P_*$ forbids no mass.
\end{enumerate}
\end{theorem}

The rotor has no continuous lift to spinors: $(e^{\theta G})^{\!\top}\eta\,e^{\theta G}=\cos2\theta\,\eta+\sin2\theta\,\eta G$, so $e^{\theta G}$ is an isometry of
$\eta$ only for $\theta\in\pi\Z$, and at $\theta=\pi$, $e^{\pi G_3}=-\id$ is implemented by $\Gamma$. This excludes a coupling of the rotor through the
spin representation and says nothing about other couplings.

\begin{theorem}[Inertia freedom at three constituents]
\label{thm:inertia3}
The real dimensions of the symmetric forms on $V_3^*$ in the Lorentz class, $\{N,G_3\}=0$, and in the duality class, $[N,G_3]=0$, that are invariant
under a given symmetry are:
\begin{center}
\begin{tabular}{@{}lcc@{}}
\toprule
symmetry & Lorentz class & duality class\\
\midrule
$\Cent(G_3)=\so{9,\C}$ & $2$ & $0$\\
$\bigoplus_k\lG^{(k)}$ & $6$ & $0$\\
the $\Pfam$-surviving algebra & $18$ & $13$\\
the $P_*$-surviving algebra & $8$ & $6$\\
none & $90$ & $81$\\
\bottomrule
\end{tabular}
\end{center}
\end{theorem}

The Lorentz-class inertia is fixed up to one complex scale under the full closure symmetry and up to three under the single-constituent Lorentz
algebras, and imposing only what a projector leaves enlarges the freedom to nine complex scales for $\Pfam$ and four for $P_*$: neither projector
selects the inertia.

\subsection{Brackets supplied by hand}

The closure above is derived; one may also supply a bracket on the frame by hand and ask what the arena's own inertias do with it. A quadratic Lie
algebra whose invariant form is $\eta$ must have a maximal compact subalgebra of exactly half its dimension, and such \emph{balanced} simple real forms
are rare: a complex simple algebra regarded as real, $\su{p,q}$ with $|p-q|=1$, or $\so{p,q}$ with $(p,q)=\bigl(\tfrac{k(k+1)}2,\tfrac{k(k-1)}2\bigr)$, and
no exceptional form. At the first rung $\R^{3,3}$ admits exactly three semisimple quadratic structures, $\so{1,3}$, $\so4$ and $\so{2,2}$, and only
$\so{1,3}$ if $\eta$ is to be proportional to the Killing form. For $n\le20$ a simple closure exists exactly at $n=1,4,5,6,7,8,12,15,16,20$.

Two observations about the arena's inertias hold across these closures. With the split-form inertia $A=m\eta$, which is an invariant form, every
solution of the Euler--Arnold equation is a one-parameter subgroup and $\dot p\equiv0$. With the positive form $A=m\,\id$ on $\so{1,3}$ the system is
$\dot\bB=0$, $\dot\bE=\tfrac2m\bE\times\bB$: the magnetic half is frozen and the electric half precesses about it. The inertias with
$G^{\!\top}A^{-1}G=-A^{-1}$ form a twelve-dimensional family of holomorphic Euler tops $\dot Z=Z\times MZ$ with $M$ complex symmetric; they are Liouville
integrable and never positive definite.

At the rungs $n=(N^2-1)/3$ with $3\nmid N$, that is $(N,n)=(2,1),(4,5),(5,8),\dots$, the compact half of the balanced form $\slr{N,\C}_\R$ is $\su N$,
which is Zeitlin's finite truncation of two-dimensional incompressible Euler flow \cite{Zeitlin91}. Giving the 't~Hooft clock--shift modes the
discrete-Laplacian inertia, numerical integration at $N=4$ and $N=5$ conserves energy and enstrophy to $10^{-12}$ and moves up to $37\%$ of the
energy and $54\%$ of the enstrophy between Laplacian shells: transfer between scales inside the frame. With either of the arena's own inertias the same
data are exact equilibria, because both $\id$ and $\eta$ restrict to multiples of the invariant form on $\su N$. What the frame lacks for a fluid is
the inertia. Finally, no closure dissipates: every Euler--Arnold flow stays on a coadjoint orbit and preserves every Casimir, which is exactly what
viscosity destroys, so no closure of the force law can produce $\nu\Delta u$ on its own; and a constant generator that uniformly decreases the split
form $C$ has, by the inertia theorem of Ostrowski and Schneider \cite{OstrowskiSchneider62}, exactly $3n$ exponentially unstable directions.
